\documentclass[11pt,fontset=fandol]{ctexart}
\usepackage[a4paper,margin=23mm]{geometry}
\usepackage{amsmath,amssymb,mathtools,bm,mathrsfs,esint}
\usepackage{graphicx,longtable,booktabs,array,calc}
\usepackage[normalem]{ulem}
\usepackage{xcolor,hyperref}
\hypersetup{hidelinks}
\setcounter{secnumdepth}{-1}
\newcounter{none}
\providecommand{\tightlist}{\setlength{\itemsep}{0pt}\setlength{\parskip}{0pt}}
\providecommand{\passthrough}[1]{#1}
\usepackage{amssymb,mathrsfs,esint,bm,newunicodechar}
\allowdisplaybreaks
\setlength{\emergencystretch}{3em}
\providecommand{\llangle}{\langle\!\langle}
\providecommand{\rrangle}{\rangle\!\rangle}
\newunicodechar{∞}{\ensuremath{\infty}}
\newunicodechar{∈}{\ensuremath{\in}}
\newunicodechar{∑}{\ensuremath{\sum}}
\newunicodechar{∫}{\ensuremath{\int}}
\newunicodechar{≤}{\ensuremath{\leq}}
\newunicodechar{≥}{\ensuremath{\geq}}
\newunicodechar{≠}{\ensuremath{\neq}}
\newunicodechar{→}{\ensuremath{\to}}
\newunicodechar{⇒}{\ensuremath{\Rightarrow}}
\newunicodechar{α}{\ensuremath{\alpha}}
\newunicodechar{β}{\ensuremath{\beta}}
\newunicodechar{γ}{\ensuremath{\gamma}}
\newunicodechar{λ}{\ensuremath{\lambda}}
\newunicodechar{φ}{\ensuremath{\varphi}}
\newunicodechar{ε}{\ensuremath{\varepsilon}}
\newunicodechar{σ}{\ensuremath{\sigma}}
\newunicodechar{θ}{\ensuremath{\theta}}
\newunicodechar{δ}{\ensuremath{\delta}}
\newunicodechar{η}{\ensuremath{\eta}}
\newunicodechar{π}{\ensuremath{\pi}}
\newunicodechar{ξ}{\ensuremath{\xi}}
\newunicodechar{Σ}{\ensuremath{\Sigma}}
\newunicodechar{Δ}{\ensuremath{\Delta}}
\newunicodechar{∂}{\ensuremath{\partial}}
\newunicodechar{∀}{\ensuremath{\forall}}
\newunicodechar{∃}{\ensuremath{\exists}}
\newunicodechar{√}{\ensuremath{\surd}}
\newunicodechar{−}{\ensuremath{-}}
\newunicodechar{μ}{\ensuremath{\mu}}
\newunicodechar{ρ}{\ensuremath{\rho}}
\newunicodechar{τ}{\ensuremath{\tau}}
\newunicodechar{ω}{\ensuremath{\omega}}
\newunicodechar{Ω}{\ensuremath{\Omega}}
\newunicodechar{χ}{\ensuremath{\chi}}
\newunicodechar{ψ}{\ensuremath{\psi}}
\newunicodechar{Γ}{\ensuremath{\Gamma}}
\newunicodechar{Φ}{\ensuremath{\Phi}}
\newunicodechar{Λ}{\ensuremath{\Lambda}}
\newunicodechar{ℝ}{\ensuremath{\mathbb{R}}}
\newunicodechar{ℤ}{\ensuremath{\mathbb{Z}}}
\newunicodechar{ℕ}{\ensuremath{\mathbb{N}}}
\newunicodechar{△}{\ensuremath{\triangle}}
\newunicodechar{⊗}{\ensuremath{\otimes}}
\begin{document}
\section{2012年杭州电子科技大学601数学分析考研真题参考答案}\label{ux5e74ux676dux5ddeux7535ux5b50ux79d1ux6280ux5927ux5b66601ux6570ux5b66ux5206ux6790ux8003ux7814ux771fux9898ux53c2ux8003ux7b54ux6848}

来源：\texttt{HDU/601数学分析/2012/2012年杭州电子科技大学601数学分析考研真题参考答案.pdf}。原
PDF 页序：1, 2, 3, 4, 5, 6, 7, 8。

\section{杭州电子科技大学2012年数学分析真题解析}\label{ux676dux5ddeux7535ux5b50ux79d1ux6280ux5927ux5b662012ux5e74ux6570ux5b66ux5206ux6790ux771fux9898ux89e3ux6790}

\begin{quote}
助手整理：按同年真题、已有本地 OCR
及可直接验证的计算恢复受损公式；以下解答经重新推导，不能视为原参考答案的逐字转录。题号与原卷对应。
\end{quote}

\subsection{一、填空题}\label{ux4e00ux586bux7a7aux9898}

\subsubsection{1.
含参变限积分}\label{ux542bux53c2ux53d8ux9650ux79efux5206}

真题及本地 OCR 均为 \(\varphi(x)=\int_x^{x^3}\arctan(y/x)\,dy\)。在
\(x\neq0\) 时令 \(y=xt\)，得 \[
\varphi(x)=x\int_1^{x^2}\arctan t\,dt,
\qquad
\boxed{\varphi'(x)=\int_1^{x^2}\arctan t\,dt+2x^2\arctan(x^2).}
\]

\subsubsection{2.
隐函数全微分}\label{ux9690ux51fdux6570ux5168ux5faeux5206}

由 \(e^{xz}+xy=z^2\)， \[
e^{xz}(z\,dx+x\,dz)+y\,dx+x\,dy=2z\,dz.
\] 因此在隐函数可解条件 \(2z-xe^{xz}\neq0\) 下， \[
\boxed{dz=\frac{(ze^{xz}+y)\,dx+x\,dy}{2z-xe^{xz}}.}
\]

\subsubsection{3. 切平面}\label{ux5207ux5e73ux9762}

对 \(z=x^2+y^2-1\)，在 \((1,1)\) 处 \(z_x=z_y=2\)，所以 \[
\boxed{2(x-1)+2(y-1)-(z-1)=0,\quad\text{即 }2x+2y-z=3.}
\]

\subsubsection{4. 定积分}\label{ux5b9aux79efux5206}

\[
\boxed{\int_1^e\frac{dx}{x\sqrt{1+\ln x}}
=\left.2\sqrt{1+\ln x}\right|_1^e=2(\sqrt2-1).}
\]

\subsubsection{5. 二阶偏导数}\label{ux4e8cux9636ux504fux5bfcux6570}

真题与本地 OCR 均为 \(u=\arctan(x/y)\)，故 \[
u_x=\frac{y}{x^2+y^2},\qquad
\boxed{u_{xx}=-\frac{2xy}{(x^2+y^2)^2}.}
\]

\subsubsection{6.
交换积分次序}\label{ux4ea4ux6362ux79efux5206ux6b21ux5e8f}

积分区域为 \(x^2+(y-1)^2\leq1\)，因此 \[
\boxed{\int_0^2dy\int_{-\sqrt{2y-y^2}}^{\sqrt{2y-y^2}}f(x,y)\,dx
=\int_{-1}^1dx\int_{1-\sqrt{1-x^2}}^{1+\sqrt{1-x^2}}f(x,y)\,dy.}
\]

\subsubsection{7.
反常积分的定义域}\label{ux53cdux5e38ux79efux5206ux7684ux5b9aux4e49ux57df}

\[
T(s)=\int_0^{+\infty}x^{s-3}e^{-x}\,dx.
\] 当 \(x\to0^+\) 时被积函数与 \(x^{s-3}\)
等价，所以零点附近收敛当且仅当
\(s-3>-1\)；无穷远处的指数衰减使其对每个实数 \(s\)
都收敛。故实参数定义域为 \[
\boxed{s\in(2,+\infty).}
\] 采用 Gamma 函数记号时，可写 \(T(s)=\Gamma(s-2)\)，该表达式在这里指
\(s>2\) 时的上述积分。

\subsubsection{8. 高阶导数}\label{ux9ad8ux9636ux5bfcux6570}

由 \(\sin^2x=(1-\cos2x)/2\)，对正整数 \(n\) 有 \[
\boxed{y^{(n)}(x)=-2^{n-1}\cos\left(2x+\frac{n\pi}{2}\right).}
\]

\subsection{二、计算题}\label{ux4e8cux8ba1ux7b97ux9898}

\subsubsection{1. 一阶导数}\label{ux4e00ux9636ux5bfcux6570}

对 \(f(x)=x^{\sin x}+x\sqrt{\sin\sqrt{x}}\)，在各项可导的定义域内部， \[
\boxed{f'(x)=x^{\sin x}\left(\cos x\ln x+\frac{\sin x}{x}\right)
+\sqrt{\sin\sqrt{x}}
+\frac{\sqrt{x}\cos\sqrt{x}}{4\sqrt{\sin\sqrt{x}}}.}
\]

\subsubsection{2. 定积分}\label{ux5b9aux79efux5206-1}

\(x^3\sin x^2\) 为奇函数，\(\cos^7(2x)\) 为偶函数，故 \[
\begin{aligned}
I&=\int_{-\pi/4}^{\pi/4}(x^3\sin x^2+\cos^7(2x))\,dx
=2\int_0^{\pi/4}\cos^7(2x)\,dx\\
&=\int_0^{\pi/2}\cos^7t\,dt
=\frac{6\cdot4\cdot2}{7\cdot5\cdot3}=\boxed{\frac{16}{35}}.
\end{aligned}
\]

\subsubsection{3. 幂级数展开}\label{ux5e42ux7ea7ux6570ux5c55ux5f00}

在 \(t=0\) 处给 \((\cos t-1)/t\)
作连续延拓。由余弦函数的幂级数，在每个有限闭区间上可逐项积分， \[
\begin{aligned}
\int_0^x\frac{\cos t-1}{t}\,dt
&=\int_0^x\sum_{n=1}^{\infty}\frac{(-1)^nt^{2n-1}}{(2n)!}\,dt\\
&=\boxed{\sum_{n=1}^{\infty}\frac{(-1)^nx^{2n}}{2n(2n)!}}.
\end{aligned}
\] 该级数的收敛半径为 \(+\infty\)，实数范围的收敛域为 \(\mathbb R\)。

\subsubsection{4. 曲线积分}\label{ux66f2ux7ebfux79efux5206}

被积形式是函数 \[
\Phi(x,y)=\frac{x^3}{3}-xy-\frac y2+\frac{\sin2y}{4}
\] 的全微分。因此从 \((0,0)\) 到 \((1,1)\) 的积分为 \[
\boxed{\Phi(1,1)-\Phi(0,0)=\frac14\sin2-\frac76.}
\]

\subsection{三、综合题}\label{ux4e09ux7efcux5408ux9898}

\subsubsection{\texorpdfstring{1. \(y=x-2\arctan x\)
的性质}{1. y=x-2\textbackslash arctan x 的性质}}\label{yx-2arctan-x-ux7684ux6027ux8d28}

\[
y'=\frac{x^2-1}{1+x^2},\qquad y''=\frac{4x}{(1+x^2)^2}.
\] 函数在 \((-\infty,-1]\)、\([1,+\infty)\) 上严格递增，在 \([-1,1]\)
上严格递减。极大值为 \(y(-1)=\pi/2-1\)，极小值为 \(y(1)=1-\pi/2\)。

在 \(x<0\) 时 \(y''<0\)，在 \(x>0\) 时 \(y''>0\)，拐点为
\((0,0)\)。用二阶导数符号说明凹凸性，可避免教材对``上凸、下凸''命名的差异。

当 \(x\to+\infty\) 时，渐近线为 \(y=x-\pi\)；当 \(x\to-\infty\)
时，渐近线为 \(y=x+\pi\)。

\subsubsection{2.
函数列一致收敛}\label{ux51fdux6570ux5217ux4e00ux81f4ux6536ux655b}

对每个 \(x\geq0\)，\(f_n(x)=xn^ke^{-nx}\to0\)。且 \[
f_n'(x)=n^k(1-nx)e^{-nx},\qquad
\sup_{x\geq0}|f_n(x)|=f_n(1/n)=e^{-1}n^{k-1}.
\] 故一致收敛当且仅当 \(\boxed{k<1}\)。\(k=1\) 时上确界恒为
\(e^{-1}\)，\(k>1\) 时上确界趋于无穷。

\subsubsection{3.
二重极限与累次极限}\label{ux4e8cux91cdux6781ux9650ux4e0eux7d2fux6b21ux6781ux9650}

\[
f(x,y)=\frac{x^2(1+x^2)-y^2(1+y^2)}{x^2+y^2}
=\frac{x^2-y^2}{x^2+y^2}+x^2-y^2.
\] 沿 \(y=0\) 时趋于 1，沿 \(x=0\) 时趋于
\(-1\)，所以原点处二重极限不存在。两种累次极限为 \[
\boxed{\lim_{x\to0}\lim_{y\to0}f(x,y)=1,\qquad
\lim_{y\to0}\lim_{x\to0}f(x,y)=-1.}
\]

\subsection{四、证明题}\label{ux56dbux8bc1ux660eux9898}

\subsubsection{\texorpdfstring{1. 极限的 \(\varepsilon\)--\(\delta\)
证明}{1. 极限的 \textbackslash varepsilon--\textbackslash delta 证明}}\label{ux6781ux9650ux7684-varepsilondelta-ux8bc1ux660e}

\[
\left|\frac{x}{x^2+1}-\frac3{10}\right|
=\frac{|3x-1||x-3|}{10(x^2+1)}.
\] 当 \(|x-3|<1\) 时，\(2<x<4\)，故 \(|3x-1|<11\) 且
\(x^2+1>5\)，上式小于 \(11|x-3|/50\)。对任意 \(\varepsilon>0\)，取 \[
\delta=\min\{1,50\varepsilon/11\}.
\] 若 \(0<|x-3|<\delta\)，则目标差小于 \(\varepsilon\)，所以极限为
\(3/10\)。

\subsubsection{2. 连续导函数}\label{ux8fdeux7eedux5bfcux51fdux6570}

题中函数为 \(f(x)=g(x)/x\)（\(x\neq0\)）、\(f(0)=g'(0)\)。因
\(g(0)=0\)，对零点附近的 \(x\) 有 \[
f(x)=\int_0^1g'(tx)\,dt.
\] 由 \(g''\) 连续，可以对参数求导，得到 \[
f'(x)=\int_0^1t\,g''(tx)\,dt.
\] 右端关于 \(x\) 连续，且 \(f'(0)=g''(0)/2\)，所以 \(f\)
有连续导函数。对非零 \(x\)，也有 \(f'(x)=[xg'(x)-g(x)]/x^2\)。

\subsubsection{3. 对数不等式}\label{ux5bf9ux6570ux4e0dux7b49ux5f0f}

对 \(u>0\)， \[
\ln(1+u)=\int_0^u\frac{dt}{1+t}>\frac{u}{1+u}.
\] 取 \(u=1/x\)，得 \(\boxed{\ln(1+1/x)>1/(1+x)}\)。

\subsubsection{4.
单调函数的黎曼可积性}\label{ux5355ux8c03ux51fdux6570ux7684ux9eceux66fcux53efux79efux6027}

单调函数在闭区间上有界。不妨设其递增，对划分
\(\Delta:a=x_0<\cdots<x_n=b\)， \[
U(f,\Delta)-L(f,\Delta)
\leq\|\Delta\|\sum_{i=1}^n[f(x_i)-f(x_{i-1})]
=\|\Delta\|[f(b)-f(a)].
\]
若端点值相同，函数为常数；否则令划分的最大间距趋于零，则上下和之差趋于零，故函数黎曼可积。递减情形同理。

\subsubsection{5.
含参反常积分的比较判别}\label{ux542bux53c2ux53cdux5e38ux79efux5206ux7684ux6bd4ux8f83ux5224ux522b}

按题目所讨论的反常积分，取被积函数在每个有限积分区间上可积。由
\(|f(x,y)|\leq F(x,y)\) 及 \(\int_a^{\infty}F(x,y)\,dx\) 关于
\(y\in[c,d]\) 一致收敛，对任意 \(\varepsilon>0\)，存在 \(M>a\)，使任意
\(v>u>M\) 及 \(y\in[c,d]\) 均满足 \[
\left|\int_u^vf(x,y)\,dx\right|
\leq\int_u^v|f(x,y)|\,dx
\leq\int_u^vF(x,y)\,dx<\varepsilon.
\] 由一致 Cauchy 准则，\(\int_a^{\infty}f(x,y)\,dx\)
一致收敛；同一估计也适用于 \(|f|\)，故该反常积分绝对收敛。
\end{document}
